%-------------------------
% Resume in LaTeX
% Author: Tanmay Kumar
% Based on: Jake's Resume (Clean, ATS-Optimized, Highly Extensible)
% License: MIT
%------------------------

\documentclass[letterpaper,10.8pt]{article}

\usepackage{latexsym}
\usepackage[empty]{fullpage}
\usepackage{titlesec}
\usepackage{marvosym}
\usepackage[usenames,dvipsnames]{color}
\usepackage{verbatim}
\usepackage{enumitem}
\usepackage[hidelinks]{hyperref}
\usepackage{fancyhdr}
\usepackage[english]{babel}
\usepackage{tabularx}
\usepackage{amsmath,amssymb}
\usepackage{xcolor}

% Set margins (0.5 in for maximum space efficiency and ATS readability)
\usepackage[margin=0.5in]{geometry}

% Hyperlink styling (sleek deep blue for interactive URLs)
\definecolor{accentblue}{RGB}{20, 65, 130}
\hypersetup{
    colorlinks=true,
    linkcolor=accentblue,
    filecolor=accentblue,      
    urlcolor=accentblue,
}

\pagestyle{fancy}
\fancyhf{} % clear all header and footer fields
\fancyfoot{}
\renewcommand{\headrulewidth}{0pt}
\renewcommand{\footrulewidth}{0pt}

% Adjust margins
\addtolength{\oddsidemargin}{-0.0in}
\addtolength{\evensidemargin}{-0.0in}
\addtolength{\textwidth}{0.0in}
\addtolength{\topmargin}{-0.2in}
\addtolength{\textheight}{0.4in}

\urlstyle{same}

\raggedbottom
\raggedright
\setlength{\tabcolsep}{0in}

% Sections formatting with clean rule
\titleformat{\section}{
  \vspace{-5pt}\scshape\raggedright\large\bfseries
}{}{0em}{}[\color{black}\titlerule \vspace{-4pt}]

% Custom commands
\newcommand{\resumeItem}[1]{
  \item\small{
    {#1 \vspace{-1.5pt}}
  }
}

\newcommand{\resumeSubheading}[4]{
  \vspace{-1pt}\item
    \begin{tabular*}{0.98\textwidth}[t]{l@{\extracolsep{\fill}}r}
      \textbf{#1} & #2 \\
      \textit{\small#3} & \textit{\small #4} \\
    \end{tabular*}\vspace{-5pt}
}

\newcommand{\resumeProjectHeading}[2]{
    \vspace{-1pt}\item
    \begin{tabular*}{0.98\textwidth}[t]{l@{\extracolsep{\fill}}r}
      #1 & #2 \\
    \end{tabular*}\vspace{-5pt}
}

\newcommand{\resumeSubHeadingListStart}{\begin{itemize}[leftmargin=0.15in, label={}]}
\newcommand{\resumeSubHeadingListEnd}{\end{itemize}}
\newcommand{\resumeItemListStart}{\begin{itemize}[leftmargin=0.15in, label={\raisebox{0.2ex}{\tiny$\bullet$}}]}
\newcommand{\resumeItemListEnd}{\end{itemize}\vspace{-4pt}}

%-------------------------------------------
%%%%%%  RESUME STARTS HERE  %%%%%%%%%%%%%%%%%
%-------------------------------------------

\begin{document}

%----------HEADING----------
\begin{center}
    {\Huge \scshape Tanmay Kumar} \\ \vspace{2pt}
    \small Indian Institute of Technology (IIT) Mandi $\cdot$ B.Tech in Civil Engineering (3rd Year) \\ \vspace{2pt}
    \small 
    \href{mailto:kumartanmay109@gmail.com}{\underline{kumartanmay109@gmail.com}} $|$ 
    \href{mailto:b24098@students.iitmandi.ac.in}{\underline{b24098@students.iitmandi.ac.in}} $|$ 
    \href{https://github.com/10tan}{\underline{github.com/10tan}} $|$ 
    \href{https://10tan.github.io}{\underline{Portfolio}} $|$
    \href{https://www.iitmandi.ac.in}{\underline{IIT Mandi}}
\end{center}
\vspace{-8pt}

%-----------EDUCATION-----------
\section{Education}
  \resumeSubHeadingListStart
    \resumeSubheading
      {Indian Institute of Technology (IIT) Mandi}{Mandi, Himachal Pradesh, India}
      {Bachelor of Technology in Civil Engineering (Current CGPA / Standing: 3rd Year)}{2022 -- 2026 (Expected)}
      \resumeItemListStart
        \resumeItem{\textbf{Focus Areas:} Computational Mechanics, Scientific Machine Learning (SciML), Stochastic Reliability Analysis, Finite Element Methods (FEM), Fluid-Structure Multiphysics, and Numerical Optimization.}
      \resumeItemListEnd
  \resumeSubHeadingListEnd
\vspace{-6pt}

%-----------TECHNICAL SKILLS-----------
\section{Technical Skills}
 \begin{itemize}[leftmargin=0.15in, label={}]
    \small{\item{
     \textbf{Languages \& Core:}{ Python (NumPy, SciPy, PyTorch, Scikit-Learn), C/C++, TypeScript, JavaScript, SQL, Bash, \LaTeX} \\
     \textbf{Scientific Computing:}{ Physics-Informed Neural Networks (PINNs, XPINNs), Fourier Neural Operators (FNO), DeepONet, Finite Element Analysis (FEA/FEM), Monte Carlo, Subset Simulation, FORM/SORM, Proper Orthogonal Decomposition (ROM)} \\
     \textbf{Engineering Standards:}{ ISO 2394 (Structural Reliability), DNV-RP-C210 (Offshore Fatigue \& Fracture), Eurocode 3 (EN 1993-1-2 Structural Fire Design), ISO 13571 (Life Safety / Tenability)} \\
     \textbf{Web \& Systems Engineering:}{ FastAPI, React 18, Vite, Three.js (WebGL 3D Rendering), Recharts, Docker, Git, CI/CD, KaTeX}
    }}
 \end{itemize}
\vspace{-6pt}

%-----------PROJECTS-----------
\section{Featured Engineering \& Scientific Projects}
    \resumeSubHeadingListStart

      % Project 1: CherenkovLens
      \resumeProjectHeading
          {\textbf{CherenkovLens} $|$ \emph{Python, Scikit-Learn, SciPy, Astrophysics, Hillas Eigensolver, Vercel}}
          {\href{https://cherenkovlens.vercel.app}{\underline{Live Demo}} $|$ \href{https://github.com/10tan/cherenkovlens}{\underline{GitHub}}}
          \resumeItemListStart
            \resumeItem{Engineered an end-to-end simulation, topological image cleaning, and machine learning pipeline for \textbf{Imaging Atmospheric Cherenkov Telescopes (IACTs)}, modeled on the 21m \textbf{MACE Observatory} in Hanle, Ladakh ($4{,}270\text{ m}$ a.s.l.).}
            \resumeItem{Simulated relativistic Extensive Air Shower (EAS) microphysics via the \textbf{Frank-Tamm equation} and barometric refractive index profiles across a $34 \times 34$ camera pixel array ($1{,}156$ PMT channels), separating electromagnetic cascades from hadronic proton backgrounds.}
            \resumeItem{Implemented dual-threshold topological tail-cuts cleaning and an analytical \textbf{Hillas eigensolver} computing statistical image moments (Length, Width, Dist, Alpha, Azwidth, Size, Concentration).}
            \resumeItem{Trained Random Forest and SVM classifiers achieving high gamma/hadron discrimination, deployed via an interactive client-side web application featuring real-time parameter visualization and KaTeX mathematical formulations.}
          \resumeItemListEnd

      % Project 2: Reliability-Lab
      \resumeProjectHeading
          {\textbf{Reliability-Lab} $|$ \emph{Python, FastAPI, React, Vite, Stochastic Reliability, Copulas, ISO 2394, DNV-RP-C210}}
          {\href{https://reliability-lab-zeta.vercel.app/}{\underline{Live Demo}} $|$ \href{https://github.com/10tan/reliability-lab}{\underline{GitHub}}}
          \resumeItemListStart
            \resumeItem{Developed an enterprise-grade structural reliability and uncertainty quantification platform compliant with \textbf{ISO 2394} and \textbf{DNV-RP-C210} probabilistic safety assessment standards.}
            \resumeItem{Formulated First- and Second-Order Reliability Methods (\textbf{FORM/SORM}) with Hasofer-Lind optimization and Rackwitz-Fiessler normal-tail transformations for nonlinear limit state functions.}
            \resumeItem{Deployed advanced rare-event simulation algorithms (\textbf{Subset Simulation, Importance Sampling, Directional Simulation}) estimating extreme failure probabilities ($P_f \in [10^{-9}, 10^{-2}]$) with low coefficient of variation.}
            \resumeItem{Modeled heavy-tailed multivariate dependencies using Archimedean copulas (Clayton, Gumbel, Frank) and Gaussian/Student-$t$ copulas; executed \textbf{Sobol variance-based global sensitivity analysis} and minimal cut-set system reliability.}
            \resumeItem{Architected a responsive full-stack platform (FastAPI + React/Vite) with real-time SVG limit-state contour plotting and automated numerical provenance verification.}
          \resumeItemListEnd

      % Project 3: firetwin
      \resumeProjectHeading
          {\textbf{firetwin} $|$ \emph{Python 3.12, PyTorch, FastAPI, React 18, TypeScript, FNO, LES CFD, Eurocode 3}}
          {\href{https://firetwin.vercel.app}{\underline{Live Demo}} $|$ \href{https://github.com/10tan/firetwin}{\underline{GitHub}}}
          \resumeItemListStart
            \resumeItem{Created a coupled multiphysics digital twin unifying Incompressible Boussinesq Navier-Stokes hydrodynamics, Smagorinsky-Lilly \textbf{Large Eddy Simulation (LES)}, and P1 radiation transport for compartment fire hazards.}
            \resumeItem{Trained a 2D \textbf{Fourier Neural Operator (FNO)} surrogate enabling continuous velocity and temperature field inference in \textbf{$<5\text{ ms}$} ($>100\times$ speedup over numerical CFD solvers) for real-time safety exploration.}
            \resumeItem{Coupled compartment gas temperatures to structural steel members via \textbf{Eurocode 3 (EN 1993-1-2)} lumped-capacitance formulations, evaluating thermal yield degradation $k_y(T)$ and elastic modulus reduction $k_E(T)$.}
            \resumeItem{Integrated \textbf{Helbing Social Force} pedestrian evacuation dynamics with ISO 13571 Fractional Effective Dose (FED) toxic incapacitation models to quantify Available Safe Egress Time (ASET) vs. Required Safe Egress Time (RSET).}
            \resumeItem{Engineered an automated serverless engine producing publication-grade 3-page compliance PDF reports complete with 6 embedded telemetry charts and Eurocode fire resistance matrices.}
          \resumeItemListEnd

      % Project 4: Neural PDE Bench
      \resumeProjectHeading
          {\textbf{Neural PDE Bench} $|$ \emph{PyTorch, Python, FastAPI, PINNs, FNO, DeepONet, Bayesian UQ, Diffusion Models}}
          {\href{https://neural-pde-bench.vercel.app}{\underline{Live Demo}} $|$ \href{https://github.com/10tan/neural-pde-bench}{\underline{GitHub}}}
          \resumeItemListStart
            \resumeItem{Built a comprehensive scientific machine learning platform integrating \textbf{Physics-Informed Neural Networks (PINNs)} and operator learning for forward and inverse partial differential equations.}
            \resumeItem{Implemented Residual-Based Adaptive Refinement (\textbf{RAR}) and multi-subdomain \textbf{XPINN decomposition} with interface flux-matching penalties to resolve stiff gradient fronts and multiscale dynamics.}
            \resumeItem{Benchmarked operator learning architectures including \textbf{Fourier Neural Operators (FNO-1D/2D)}, \textbf{DeepONet} (branch-trunk networks), and Geometry Graph Neural Operators (GNO) across continuum mechanics benchmarks.}
            \resumeItem{Conducted Bayesian Uncertainty Quantification via \textbf{Hamiltonian Monte Carlo (HMC)} and Variational MC-Dropout; deployed score-based diffusion generative models for inverse spatial parameter reconstruction $E(x)$.}
            \resumeItem{Constructed an interactive React/TypeScript analytical suite with real-time Recharts dashboards and a containerized FastAPI backend with automated pytest CI pipelines.}
          \resumeItemListEnd

      % Project 5: Nonlinear FEA Studio
      \resumeProjectHeading
          {\textbf{Nonlinear FEA Studio} $|$ \emph{Python, SciPy, Three.js, WebGL, React, Arc-Length Continuation, ROM, DNV Standards}}
          {\href{https://nonlinear-fea-studio.vercel.app}{\underline{Live Demo}} $|$ \href{https://github.com/10tan/nonlinear-fea-studio}{\underline{GitHub}}}
          \resumeItemListStart
            \resumeItem{Engineered an advanced computational structural mechanics platform supporting geometrically and materially nonlinear continuum regimes, multiscale homogenization, and structural optimization.}
            \resumeItem{Formulated Total and Updated Lagrangian kinematics (Green-Lagrange strains, 2nd Piola-Kirchhoff stresses) coupled with \textbf{Crisfield/Riks cylindrical arc-length continuation} to trace snap-through and snap-back buckling paths.}
            \resumeItem{Derived return-mapping $J_2$ plasticity with isotropic/kinematic hardening and \textbf{Lemaitre continuum damage mechanics} ($\dot{D}$), accompanied by Zienkiewicz-Zhu adaptive mesh refinement (AMR).}
            \resumeItem{Integrated \textbf{SIMP topology optimization} with MMA optimizer and manufacturing overhang constraints; computed asymptotic microstructural homogenization ($FE^2$) with Voigt/Reuss/Hashin-Shtrikman bounds.}
            \resumeItem{Developed snapshot \textbf{POD-Galerkin reduced-order models (ROM)} for real-time load sliders and built an interactive Three.js 3D WebGL CAD interface visualizing stress tensors, plastic strain, and ASTM E1049 rainflow fatigue life.}
          \resumeItemListEnd

    \resumeSubHeadingListEnd
\vspace{-6pt}

%-----------RELEVANT COURSEWORK & HONORS-----------
\section{Relevant Coursework \& Academic Highlights}
 \begin{itemize}[leftmargin=0.15in, label={}]
    \small{\item{
     \textbf{Computational \& Mathematical Foundations:}{ Numerical Methods in Engineering, Finite Element Analysis, Probability \& Stochastic Processes, Linear Algebra, Multivariable Calculus, Optimization Techniques, Differential Equations} \\
     \textbf{Structural \& Physical Mechanics:}{ Structural Analysis, Solid Mechanics, Continuum Mechanics, Fluid Mechanics, Structural Dynamics, Mechanics of Materials, Probabilistic Risk Assessment} \\
     \textbf{Research Interests:}{ Scientific Machine Learning (Physics-Informed Deep Learning \& Operator Learning), Structural Health Monitoring, Computational Fracture Mechanics, Multi-Hazard Digital Twins, and Offshore Engineering Reliability}
    }}
 \end{itemize}

\end{document}
